\begin{textsample}{POS Dim 1 – 2008 – Score 24.00 – 2008/t000180.txt}  \label{ex:f1_pos_008}
I thought I ' d start with telling you or showing you \textbf{the} people who started [ Jet Propulsion Lab ].
When they were a bunch of kids, they were kind of very imaginative, very adventurous, as they were trying at Caltech to mix chemicals and see which one blows up more.
Well, I don ' t recommend that you try to do that now.
Naturally, they blew up a shack, and Caltech, well, then, hey, you go to \textbf{the} Arroyo and really do all your tests in there.
So, that ' s what we call our first five employees during \textbf{the} tea break, you know, in here.
As I said, they were adventurous people.
As a matter of fact, one of them, who was, kind of, part of a cult which was not too far from here on Orange Grove, and unfortunately he blew up himself because he kept mixing chemicals and trying to figure out which ones were \textbf{the} best chemicals.
So, that gives you a kind of flavor of \textbf{the} kind of people we have there.
We try to avoid blowing ourselves up.
This one I thought I ' d show you.
Guess which one is a JPL employee in \textbf{the} heart of this crowd.
I tried to come like him this morning, but as I walked out, then it was too cold, and I said, I ' d better put my shirt back on.
But more importantly, \textbf{the} reason I wanted to show this picture : look where \textbf{the} other people are looking, and look where he is looking.
Wherever anybody else looks, look somewhere else, and go do something different, you know, and doing that.
And that ' s kind of what has been \textbf{the} spirit of what we are doing.
And I want to tell you a quote from Ralph Emerson that one of my colleagues, you know, put on my wall in my office, and it says,"Do not go where \textbf{the} path may lead.
Go instead where there is no path, and leave a trail."And that ' s my recommendation to all of you : look what everybody is doing, what they are doing ; go do something completely different.
Don ' t try to improve a little bit on what somebody else is doing, because that doesn ' t get you very far.
In our early days we used to work a lot on rockets, but we also used to have a lot of parties, you know.
As you can see, one of our parties, you know, a few years ago.
But then a big difference happened about 50 years ago, after Sputnik was launched.
We launched \textbf{the} first American satellite, and that ' s \textbf{the} one you see on \textbf{the} \textbf{left} in there.
And here we made 180 degrees change : we changed from a rocket house to be an exploration house.
And that was done over a period of a couple of years, and now we are \textbf{the} leading organization, you know, exploring space on all of your behalf.
But even when we did that, we had to remind ourselves, sometimes there are setbacks.
So you see, on \textbf{the} bottom, that rocket was supposed to go upward ; somehow it ended going sideways.
So that ' s what we call \textbf{the} misguided missile.
But then also, just to celebrate that, we started an event at JPL for"Miss Guided Missile."So, we used to have a celebration every year and select — there used to be competition and parades and so on.
It ' s not very appropriate to do it now.
Some people tell me to do it ; I think, well, that ' s not really proper, you know, these days.
So, we do something a little bit more serious.
And that ' s what you see in \textbf{the} last Rose Bowl, you know, when we entered one of \textbf{the} floats.
That ' s more on \textbf{the} play side.
And on \textbf{the} right side, that ' s \textbf{the} Rover just before we finished its testing to take it to \textbf{the} Cape to launch it.
These are \textbf{the} Rovers up here that you have on Mars now.
So that kind of tells you about, kind of, \textbf{the} fun things, you know, and \textbf{the} serious things that we try to do.
But I said I ' m going to show you a short \textbf{clip} of one of our employees to kind of give you an idea about some of \textbf{the} talent that we have.
Video : Morgan Hendry : Beware of Safety is an instrumental \textbf{rock} band.
It branches on more \textbf{the} experimental side.
There ' s \textbf{the} improvisational side of jazz.
There ' s \textbf{the} heavy-hitting sound of \textbf{rock}.
Being able to treat sound as an instrument, and be able to dig for more \textbf{abstract} sounds and things to play live, mixing electronics and acoustics. \textbf{The} music ' s half of me, but \textbf{the} other half — I landed probably \textbf{the} best gig of all.
I work for \textbf{the} Jet Propulsion Lab.
I ' m building \textbf{the} next Mars Rover.
Some of \textbf{the} most brilliant engineers I know are \textbf{the} ones who have that sort of artistic quality about them.
You ' ve got to do what you want to do.
And anyone who tells you you can ' t, you don ' t listen to them.
Maybe they ' re right - I doubt it.
Tell them where to put it, and then just do what you want to do.
I ' m Morgan Hendry.
I am NASA.
Charles Elachi : Now, moving from \textbf{the} play stuff to \textbf{the} serious stuff, always people ask, why do we explore?
Why are we doing all of these missions and why are we exploring them?
Well, \textbf{the} way I think about it is fairly simple.
Somehow, 13 billion years ago there was a Big Bang, and you ' ve heard a little bit about, you know, \textbf{the} origin of \textbf{the} \textbf{universe}.
But somehow what strikes everybody ' s imagination — or lots of people ' s imagination — somehow from that original Big Bang we have this beautiful world that we live in today.
You look outside : you have all that beauty that you see, all that life that you see around you, and here we have intelligent people like you and I who are having a conversation here.
All that started from that Big Bang.
So, \textbf{the} question is : How did that happen?
How did that \textbf{evolve}?
How did \textbf{the} \textbf{universe} form?
How did \textbf{the} galaxies form?
How did \textbf{the} \textbf{planets} form?
Why is there a \textbf{planet} on which there is life which have \textbf{evolved}?
Is that very common?
Is there life on every \textbf{planet} that you can see around \textbf{the} stars?
So we literally are all made out of stardust.
We started from those stars ; we are made of stardust.
So, next time you are really depressed, look in \textbf{the} mirror and you can look and say, hi, I ' m looking at a star here.
You can skip \textbf{the} dust part.
But literally, we are all made of stardust.
So, what we are trying to do in our exploration is effectively write \textbf{the} book of how things have came about as they are today.
And one of \textbf{the} first, or \textbf{the} easiest, places we can go and explore that is to go towards Mars.
And \textbf{the} reason Mars takes particular attention : it ' s not very far from us.
You know, it ' ll take us only six months to get there.
Six to nine months at \textbf{the} right time of \textbf{the} year.
It ' s a \textbf{planet} somewhat similar to Earth.
It ' s a little bit smaller, but \textbf{the} land mass on Mars is about \textbf{the} same as \textbf{the} land mass on Earth, you know, if you don ' t take \textbf{the} oceans into account.
It has polar caps.
It has an atmosphere somewhat thinner than ours, so it has weather.
So, it ' s very similar to some extent, and you can see some of \textbf{the} features on it, like \textbf{the} Grand Canyon on Mars, or what we call \textbf{the} Grand Canyon on Mars.
It is like \textbf{the} Grand Canyon on Earth, except a hell of a lot larger.
So it ' s about \textbf{the} size, you know, of \textbf{the} United States.
It has volcanoes on it.
And that ' s Mount Olympus on Mars, which is a kind of huge volcanic shield on that \textbf{planet}.
And if you look at \textbf{the} height of it and you compare it to Mount Everest, you see, it ' ll give you an idea of how large that Mount Olympus, you know, is, \textbf{relative} to Mount Everest.
So, it basically \textbf{dwarfs}, you know, Mount Everest here on Earth.
So, that gives you an idea of \textbf{the} tectonic events or volcanic events which have happened on that \textbf{planet}.
Recently from one of our satellites, this shows that it ' s Earth-like — we caught a landslide \textbf{occurring} as it was happening.
So it is a dynamic \textbf{planet}, and activity is going on as we speak today.
And these Rovers, people wonder now, what are they doing today, so I thought I would show you a little bit what they are doing.
This is one very large crater.
Geologists love craters, because craters are like digging a big hole in \textbf{the} ground without really working at it, and you can see what ' s below \textbf{the} \textbf{surface}.
So, this is called Victoria Crater, which is about a few football fields in size.
And if you look at \textbf{the} top \textbf{left}, you see a little teeny dark dot.
This picture was taken from an orbiting satellite.
If I zoom on it, you can see : that ' s \textbf{the} Rover on \textbf{the} \textbf{surface}.
So, that was taken from orbit ; we had \textbf{the} camera zoom on \textbf{the} \textbf{surface}, and we actually saw \textbf{the} Rover on \textbf{the} \textbf{surface}.
And we actually used \textbf{the} combination of \textbf{the} satellite images and \textbf{the} Rover to actually conduct science, because we can observe large areas and then you can get those Rovers to move around and basically go to a certain location.
So, specifically what we are doing now is that Rover is going down in that crater.
As I told you, geologists love craters.
And \textbf{the} reason is, many of you went to \textbf{the} Grand Canyon, and you see in \textbf{the} wall of \textbf{the} Grand Canyon, you see these layers.
And what these layers — that ' s what \textbf{the} \textbf{surface} used to be a million years ago, 10 million years ago, 100 million years ago, and you get deposits on top of them.
So if you can read \textbf{the} layers it ' s like reading your book, and you can learn \textbf{the} history of what happened in \textbf{the} past in that location.
So what you are seeing here are \textbf{the} layers on \textbf{the} wall of that crater, and \textbf{the} Rover is going down now, measuring, you know, \textbf{the} properties and analyzing \textbf{the} \textbf{rocks} as it ' s going down, you know, that canyon.
Now, it ' s kind of a little bit of a challenge driving down a \textbf{slope} like this.
If you were there you wouldn ' t do it yourself.
But we really made sure we tested those Rovers before we got them down — or that Rover — and made sure that it ' s all working well.
Now, when I came last time, shortly after \textbf{the} landing — I think it was, like, a hundred days after \textbf{the} landing — I told you I was surprised that those Rovers are lasting even a hundred days.
Well, here we are four years later, and they ' re still working.
Now you say, Charles, you are really lying to us, and so on, but that ' s not true.
We really believed they were going to last 90 days or 100 days, because they are solar powered, and Mars is a dusty \textbf{planet}, so we expected \textbf{the} dust would start accumulating on \textbf{the} \textbf{surface}, and after a while we wouldn ' t have enough power, you know, to keep them warm.
Well, I always say it ' s important that you are smart, but every once in a while it ' s good to be lucky.
And that ' s what we found out.
It turned out that every once in a while there are dust devils which come by on Mars, as you are seeing here, and when \textbf{the} dust devil comes over \textbf{the} Rover, it just cleans it up.
It is like a brand new car that you have, and that ' s literally why they have lasted so long.
And now we designed them reasonably well, but that ' s exactly why they are lasting that long and still providing all \textbf{the} science data.
Now, \textbf{the} two Rovers, each one of them is, kind of, getting old.
You know, one of them, one of \textbf{the} wheels is stuck, is not working, one of \textbf{the} front wheels, so what we are doing, we are driving it \textbf{backwards}.
And \textbf{the} other one has arthritis of \textbf{the} shoulder joint, you know, it ' s not working very well, so it ' s walking like this, and we can move \textbf{the} arm, you know, that way.
But still they are producing a lot of scientific data.
Now, during that whole period, a number of people got excited, you know, outside \textbf{the} science community about these Rovers, so I thought I ' d show you a video just to give you a reflection about how these Rovers are being viewed by people other than \textbf{the} science community.
So let me go on \textbf{the} next short video.
By \textbf{the} way, this video is pretty accurate of how \textbf{the} landing took place, you know, about four years ago.
Video : \textbf{Okay}, we have parachute aligned. \textbf{Okay}, deploy \textbf{the} airbags.
Open.
Camera.
We have a picture right now.
Yeah!
CE : That ' s about what happened in \textbf{the} Houston operation room.
It ' s exactly like this.
Video : Now, if there is life, \textbf{the} Dutch will find it.
What is he doing?
What is that?
CE : Not too bad.
So anyway, let me continue on showing you a little bit about \textbf{the} beauty of that \textbf{planet}.
As I said earlier, it looked very much like Earth, so you see sand dunes.
It looks like I could have told you these are pictures taken from \textbf{the} Sahara Desert or somewhere, and you ' d have believed me, but these are pictures taken from Mars.
But one area which is particularly intriguing for us is \textbf{the} northern region, you know, of Mars, close to \textbf{the} North Pole, because we see ice caps, and we see \textbf{the} ice caps shrinking and expanding, so it ' s very much like you have in northern Canada.
And we wanted to find out — and we see all kinds of glacial features on it.
So, we wanted to find out, actually, what is that ice made of, and could that have embedded in it some \textbf{organic}, you know, material.
So we have a spacecraft which is heading towards Mars, called Phoenix, and that spacecraft will land 17 days, seven hours and 20 seconds from now, so you can adjust your watch.
So it ' s on May 25 around just before five o ' clock our time here on \textbf{the} West Coast, actually we will be landing on another \textbf{planet}.
And as you can see, this is a picture of \textbf{the} spacecraft put on Mars, but I thought that just in case you ' re going to miss that show, you know, in 17 days, I ' ll show you, kind of, a little bit of what ' s going to happen.
Video : That ' s what we call \textbf{the} seven minutes of terror.
So \textbf{the} plan is to dig in \textbf{the} soil and take samples that we put them in an oven and actually heat them and look what gases will come from it.
So this was launched about nine months ago.
We ' ll be coming in at 12, 000 miles per hour, and in seven minutes we have to stop and touch \textbf{the} \textbf{surface} very softly so we don ' t break that lander.
Ben Cichy : Phoenix is \textbf{the} first Mars Scout mission.
It ' s \textbf{the} first mission that ' s going to try to land near \textbf{the} North Pole of Mars, and it ' s \textbf{the} first mission that ' s actually going to try and reach out and touch water on \textbf{the} \textbf{surface} of another \textbf{planet}.
Lynn Craig : Where there tends to be water, at least on Earth, there tends to be life, and so it ' s potentially a place where life could have existed on \textbf{the} \textbf{planet} in \textbf{the} past.
Erik Bailey : \textbf{The} main purpose of EDL is to take a spacecraft that is traveling at 12, 500 miles an hour and bring it to a screeching halt in a soft way in a very short amount of time.
BC : We enter \textbf{the} Martian atmosphere.
We ' re 70 miles above \textbf{the} \textbf{surface} of Mars.
And our lander is safely tucked inside what we call an aeroshell.
EB : Looks kind of like an ice cream cone, more or less.
BC : And on \textbf{the} front of it is this heat shield, this saucer-looking thing that has about a half-inch of essentially what ' s cork on \textbf{the} front of it, which is our heat shield.
Now, this is really special cork, and this cork is what ' s going to protect us from \textbf{the} violent atmospheric entry that we ' re about to experience.
Rob Grover : Friction really starts to build up on \textbf{the} spacecraft, and we use \textbf{the} friction when it ' s flying through \textbf{the} atmosphere to our advantage to slow us down.
BC : From this point, we ' re going to decelerate from 12, 500 miles an hour down to 900 miles an hour.
EB : \textbf{The} outside can get almost as hot as \textbf{the} \textbf{surface} of \textbf{the} Sun.
RG : \textbf{The} temperature of \textbf{the} heat shield can reach 2, 600 degrees Fahrenheit.
EB : \textbf{The} inside doesn ' t get very hot.
It probably gets about room temperature.
Richard Kornfeld : There is this window of opportunity within which we can deploy \textbf{the} parachute.
EB : If you fire \textbf{the} ' chute too early, \textbf{the} parachute itself could fail. \textbf{The} fabric and \textbf{the} stitching could just pull apart.
And that would be bad.
BC : In \textbf{the} first 15 seconds after we deploy \textbf{the} parachute, we ' ll decelerate from 900 miles an hour to a relatively slow 250 miles an hour.
We no longer need \textbf{the} heat shield to protect us from \textbf{the} force of atmospheric entry, so we jettison \textbf{the} heat shield, exposing for \textbf{the} first time our lander to \textbf{the} atmosphere of Mars.
LC : After \textbf{the} heat shield has been jettisoned and \textbf{the} legs are deployed, \textbf{the} next step is to have \textbf{the} radar system begin to \textbf{detect} how far Phoenix really is from \textbf{the} ground.
BC : We ' ve lost 99 percent of our entry velocity.
So, we ' re 99 percent of \textbf{the} way to where we want to be.
But that last one percent, as it always seems to be, is \textbf{the} tricky part.
EB : Now \textbf{the} spacecraft actually has to decide when it ' s going to get rid of its parachute.
BC : We separate from \textbf{the} lander going 125 miles an hour at roughly a kilometer above \textbf{the} \textbf{surface} of Mars : 3, 200 feet.
That ' s like taking two Empire State Buildings and \textbf{stacking} them on top of one another.
EB : That ' s when we separate from \textbf{the} back shell, and we ' re now in free-fall.
It ' s a very scary moment ; a lot has to happen in a very short amount of time.
LC : So it ' s in a free-fall, but it ' s also trying to use all of its actuators to make sure that it ' s in \textbf{the} right position to land.
EB : And then it has to light up its engines, right itself, and then slowly slow itself down and touch down on \textbf{the} ground safely.
BC : Earth and Mars are so far apart that it takes over ten minutes for a signal from Mars to get to Earth.
And EDL itself is all over in a matter of seven minutes.
So by \textbf{the} time you even hear from \textbf{the} lander that EDL has started it ' ll already be over.
EB : We have to build large amounts of autonomy into \textbf{the} spacecraft so that it can land itself safely.
BC : EDL is this immense, technically challenging problem.
It ' s about getting a spacecraft that ' s hurtling through deep space and using all this bag of tricks to somehow figure out how to get it down to \textbf{the} \textbf{surface} of Mars at zero miles an hour.
It ' s this immensely exciting and challenging problem.
CE : Hopefully it all will happen \textbf{the} way you saw it in here.
So it will be a very tense moment, you know, as we are watching that spacecraft landing on another \textbf{planet}.
So now let me talk about \textbf{the} next things that we are doing.
So we are in \textbf{the} process, as we speak, of actually designing \textbf{the} next Rover that we are going to be sending to Mars.
So I thought I would go a little bit and tell you, kind of, \textbf{the} steps we go through.
It ' s very similar to what you do when you design your product.
As you saw a little bit earlier, when we were doing \textbf{the} Phoenix one, we have to take into account \textbf{the} heat that we are going to be facing.
So we have to study all kinds of different materials, \textbf{the} shape that we want to do.
In general we don ' t try to please \textbf{the} customer here.
What we want to do is to make sure we have an effective, you know, an efficient kind of machine.
First we start by we want to have our employees to be as imaginative as they can.
And we really love being close to \textbf{the} art center, because we have, as a matter of fact, one of \textbf{the} alumni from \textbf{the} art center, Eric Nyquist, had put a series of displays, far-out displays, you know, in our what we call mission design or spacecraft design room, just to get people to think wildly about things.
We have a bunch of Legos.
So, as I said, this is a playground for adults, where they sit down and try to play with different shapes and different designs.
Then we get a little bit more serious, so we have what we call our CAD / CAMs and all \textbf{the} engineers who are involved, or scientists who are involved, who know about thermal properties, know about design, know about atmospheric interaction, parachutes, all of these things, which they work in a team effort and actually design a spacecraft in a \textbf{computer} to some extent, so to see, does that meet \textbf{the} requirement that we need.
On \textbf{the} right, also, we have to take into account \textbf{the} environment of \textbf{the} \textbf{planet} where we are going.
If you are going to Jupiter, you have a very high-radiation, you know, environment.
It ' s about \textbf{the} same radiation environment close by Jupiter as inside a nuclear reactor.
So just imagine : you take your P.
C. and throw it into a nuclear reactor and it still has to work.
So these are kind of some of \textbf{the} little challenges, you know, that we have to face.
If we are doing entry, we have to do tests of parachutes.
You saw in \textbf{the} video a parachute breaking.
That would be a bad day, you know, if that happened, so we have to test, because we are deploying this parachute at supersonic speeds.
We are coming at extremely high speeds, and we are deploying them to slow us down.
So we have to do all kinds of tests.
To give you an idea of \textbf{the} size, you know, of that parachute \textbf{relative} to \textbf{the} people standing there.
Next step, we go and actually build some kind of test models and actually test them, you know, in \textbf{the} lab at JPL, in what we call our Mars Yard.
We kick them, we hit them, we drop them, just to make sure we understand how, where would they break.
And then we back off, you know, from that point.
And then we actually do \textbf{the} actual building and \textbf{the} flight.
And this next Rover that we ' re flying is about \textbf{the} size of a car.
That big shield that you see outside, that ' s a heat shield which is going to protect it.
And that will be basically built over \textbf{the} next year, and it will be launched June a year from now.
Now, in that case, because it was a very big Rover, we couldn ' t use airbags.
And I know many of you, kind of, last time afterwards said well, that was a cool thing to have — those airbags.
Unfortunately this Rover is, like, ten times \textbf{the} size of \textbf{the}, you know, mass-wise, of \textbf{the} other Rover, or three times \textbf{the} mass.
So we can ' t use airbags.
So we have to come up with another ingenious idea of how do we land it.
And we didn ' t want to take it propulsively all \textbf{the} way to \textbf{the} \textbf{surface} because we didn ' t want to contaminate \textbf{the} \textbf{surface} ; we wanted \textbf{the} Rover to immediately land on its legs.
So we came up with this ingenious idea, which is used here on Earth for helicopters.
Actually, \textbf{the} lander will come down to about 100 feet and hover above that \textbf{surface} for 100 feet, and then we have a sky crane which will take that Rover and land it down on \textbf{the} \textbf{surface}.
Hopefully it all will work, you know, it will work that way.
And that Rover will be more kind of like a chemist.
What we are going to be doing with that Rover as it drives around, it ' s going to go and analyze \textbf{the} chemical composition of \textbf{rocks}.
So it will have an arm which will take samples, put them in an oven, crush and analyze them.
But also, if there is something that we cannot reach because it is too high on a cliff, we have a little laser system which will actually zap \textbf{the} \textbf{rock}, evaporate some of it, and actually analyze what ' s coming from that \textbf{rock}.
So it ' s a little bit like"Star Wars,"you know, but it ' s real.
It ' s real stuff.
And also to help you, to help \textbf{the} community so you can do ads on that Rover, we are going to train that Rover to actually in addition to do this, to actually serve cocktails, you know, also on Mars.
So that ' s kind of giving you an idea of \textbf{the} kind of, you know, fun things we are doing on Mars.
I thought I ' d go to"\textbf{The} Lord of \textbf{the} Rings"now and show you some of \textbf{the} things we have there.
Now,"\textbf{The} Lord of \textbf{the} Rings"has two things played through it.
One, it ' s a very attractive \textbf{planet} — it just has \textbf{the} beauty of \textbf{the} rings and so on.
But for scientists, also \textbf{the} rings have a special meaning, because we believe they represent, on a small scale, how \textbf{the} Solar System actually formed.
Some of \textbf{the} scientists believe that \textbf{the} way \textbf{the} Solar System formed, that \textbf{the} Sun when it collapsed and actually created \textbf{the} Sun, a lot of \textbf{the} dust around it created rings and then \textbf{the} particles in those rings accumulated together, and they formed bigger \textbf{rocks}, and then that ' s how \textbf{the} \textbf{planets}, you know, were formed.
So, \textbf{the} idea is, by watching Saturn we ' re actually watching our solar system in real time being formed on a smaller scale, so it ' s like a test bed for it.
So, let me show you a little bit on what that Saturnian system looks like.
First, I ' m going to fly you over \textbf{the} rings.
By \textbf{the} way, all of this is real stuff.
This is not animation or anything like this.
This is actually taken from \textbf{the} satellite that we have in orbit around Saturn, \textbf{the} Cassini.
And you see \textbf{the} amount of detail that is in those rings, which are \textbf{the} particles.
Some of them are agglomerating together to form larger particles.
So that ' s why you have these gaps, is because a small satellite, you know, is being formed in that location.
Now, you think that those rings are very large \textbf{objects}.
Yes, they are very large in one \textbf{dimension} ; in \textbf{the} other \textbf{dimension} they are paper thin.
Very, very thin.
What you are seeing here is \textbf{the} shadow of \textbf{the} ring on Saturn itself.
And that ' s one of \textbf{the} satellites which was actually formed on that one.
So, think about it as a paper-thin, huge area of many hundreds of thousands of miles, which is rotating.
And we have a wide variety of kind of satellites which will form, each one looking very different and very odd, and that keeps scientists busy for tens of years trying to explain this, and telling NASA we need more money so we can explain what these things look like, or why they formed that way.
Well, there were two satellites which were particularly interesting.
One of them is called Enceladus.
It ' s a satellite which was all made of ice, and we measured it from orbit.
Made of ice.
But there was something bizarre about it.
If you look at these stripes in here, what we call \textbf{tiger} stripes, when we flew over them, all of a sudden we saw an increase in \textbf{the} temperature, which said that those stripes are warmer than \textbf{the} rest of \textbf{the} \textbf{planet}.
So as we flew by away from it, we looked back.
And guess what?
We saw geysers coming out.
So this is a Yellowstone, you know, of Saturn.
We are seeing geysers of ice which are coming out of that \textbf{planet}, which indicate that most likely there is an ocean, you know, below \textbf{the} \textbf{surface}.
And somehow, through some dynamic effect, we ' re having these geysers which are being, you know, emitted from it.
And \textbf{the} reason I showed \textbf{the} little arrow there, I think that should say 30 miles, we decided a few months ago to actually fly \textbf{the} spacecraft through \textbf{the} plume of that geyser so we can actually measure \textbf{the} material that it is made of.
That was [ \textbf{unclear} ] also — you know, because we were worried about \textbf{the} risk of it, but it worked pretty well.
We flew at \textbf{the} top of it, and we found that there is a fair amount of \textbf{organic} material which is being emitted in combination with \textbf{the} ice.
And over \textbf{the} next few years, as we keep orbiting, you know, Saturn, we are planning to get closer and closer down to \textbf{the} \textbf{surface} and make more accurate measurements.
Now, another satellite also attracted a lot of attention, and that ' s Titan.
And \textbf{the} reason Titan is particularly interesting, it ' s a satellite bigger than our moon, and it has an atmosphere.
And that atmosphere is very — as dense as our own atmosphere.
So if you were on Titan, you would feel \textbf{the} same pressure that you feel in here.
Except it ' s a lot colder, and that atmosphere is heavily made of methane.
Now, methane gets people all excited, because it ' s \textbf{organic} material, so immediately people start thinking, could life have \textbf{evolved} in that location, when you have a lot of \textbf{organic} material.
So people believe now that Titan is most likely what we call a pre-biotic \textbf{planet}, because it ' s so cold \textbf{organic} material did not get to \textbf{the} stage of becoming biological material, and therefore life could have \textbf{evolved} on it.
So it could be Earth, frozen three billion years ago before life actually started on it.
So that ' s getting a lot of interest, and to show you some example of what we did in there, we actually dropped a probe, which was developed by our colleagues in Europe, we dropped a probe as we were orbiting Saturn.
We dropped a probe in \textbf{the} atmosphere of Titan.
And this is a picture of an area as we were coming down.
Just looked like \textbf{the} coast of California for me.
You see \textbf{the} rivers which are coming along \textbf{the} coast, and you see that white area which looks like Catalina Island, and that looks like an ocean.
And then with an instrument we have on board, a radar instrument, we found there are lakes like \textbf{the} Great Lakes in here, so it looks very much like Earth.
It looks like there are rivers on it, there are oceans or lakes, we know there are clouds.
We think it ' s raining also on it.
So it ' s very much like \textbf{the} cycle on Earth except because it ' s so cold, it could not be water, you know, because water would have frozen.
What it turned out, that all that we are seeing, all this liquid, [ is made of ] hydrocarbon and ethane and methane, similar to what you put in your car.
So here we have a cycle of a \textbf{planet} which is like our Earth, but is all made of ethane and methane and \textbf{organic} material.
So if you were on Mars — sorry, on Titan, you don ' t have to worry about four-dollar gasoline.
You just drive to \textbf{the} nearest lake, stick your hose in it, and you ' ve got your car filled up.
On \textbf{the} other hand, if you light a match \textbf{the} whole \textbf{planet} will blow up.
So in closing, I said I want to close by a couple of pictures.
And just to kind of put us in perspective, this is a picture of Saturn taken with a spacecraft from behind Saturn, looking towards \textbf{the} Sun. \textbf{The} Sun is behind Saturn, so we see what we call"forward scattering,"so it highlights all \textbf{the} rings.
And I ' m going to zoom.
There is a — I ' m not sure you can see it very well, but on \textbf{the} top \textbf{left}, around 10 o ' clock, there is a little teeny dot, and that ' s Earth.
You barely can see ourselves.
So what I did, I thought I ' d zoom on it.
So as you zoom in, you know, you can see Earth, you know, just in \textbf{the} middle here.
So we zoomed all \textbf{the} way on \textbf{the} art center.
So thank you very much.

% matched lemmas: abstract, backwards, clip, computer, detect, dimension, dwarf, evolve, left, object, occur, okay, organic, planet, relative, rock, slope, stack, surface, the, tiger, unclear, universe
\end{textsample}
